\begin{proofbox}
	\textbf{\textcolor{CProof}{思路提示}}

	利用正弦定理及外角平分线补角性质，通过比例等式变换得到结论。

	\medskip
	\textbf{\textcolor{CProof}{正式推导}}
	\begin{description}[style=nextline, leftmargin=2.8em, labelsep=0.5em, itemsep=0.55em]
		\item[\textbf{步骤一（设角得补角）：}]
			\[
				\angle BAD = \beta,\quad \angle CAD = 180^\circ - \beta.
			\]
			\par\textbf{理由：} 由$AD$是$\angle A$的外角平分线可知$\angle BAD$与$\angle CAD$互补。

		\item[\textbf{步骤二（ABD正弦定理）：}]
			\[
				\frac{BD}{\sin\beta} = \frac{AB}{\sin\angle ADB}.
			\]
			\par\textbf{理由：} 在$\triangle ABD$中应用正弦定理，将$\sin\angle BAD$替换为$\sin\beta$。

		\item[\textbf{步骤三（ACD正弦定理）：}]
			\[
				\frac{DC}{\sin(180^\circ-\beta)} = \frac{AC}{\sin\angle ADC},\quad \frac{DC}{\sin\beta} = \frac{AC}{\sin\angle ADC}.
			\]
			\par\textbf{理由：} 在$\triangle ACD$中应用正弦定理，并由$\sin(180^\circ-\beta)=\sin\beta$化简。

		\item[\textbf{步骤四（共线正弦相等）：}]
			\[
				\angle ADB+\angle ADC=180^\circ,\quad \sin\angle ADB=\sin\angle ADC.
			\]
			\par\textbf{理由：} 因为$B,D,C$共线，所以$\angle ADB$与$\angle ADC$互补，从而正弦相等。

		\item[\textbf{步骤五（相除得比例）：}]
			\[
				\frac{BD/\sin\beta}{DC/\sin\beta} = \frac{AB/\sin\angle ADB}{AC/\sin\angle ADC} \ \Longrightarrow\ \frac{BD}{DC} = \frac{AB}{AC}.
			\]
			\par\textbf{理由：} 将前两等式相除，约去$\sin\beta$和$\sin\angle ADB=\sin\angle ADC$。

	\end{description}

	\medskip
	\textbf{\textcolor{CProof}{结论回扣}}

	\textbf{由此得到外角平分线定理的结论：} $\displaystyle \frac{BD}{DC} = \frac{AB}{AC}$。
\end{proofbox}
